% !PRINTED 156
\weberpage{156}
%<<CONT
von bestimmter Dimension eine (unendliche) \gesperrt{Gruppe} (\S.\,1). Wenn wir aus dieser Gesammtheit irgend eine Menge herausheben, die so in sich abgeschlossen ist, dass irgend zwei ihrer Elemente durch Composition ein Element derselben Menge ergeben, so bildet diese Menge gleichfalls eine \gesperrt{Gruppe}, die endlich oder unendlich sein kann.

Bedeutet $L$ eine feste Substitution, so kann man aus jeder Substitution $A$ von derselben Dimension eine Substitution

\weq{21}{A' = L^{-1}AL}

ableiten, die die \gesperrt{Transformirte von $A$ durch $L$} heisst.

Setzen wir

\weq{22}{(x) = L(x'), \quad (y) = L(y'),}

so folgt aus (4):

\weq{23}{(y') = A'(x'),}

so dass der Uebergang zu der transformirten Substitution gleichbedeutend ist mit der gleichzeitigen Transformation beider Reihen von Veränderlichen durch $L^{-1}$.

Ist

\[
A' = L^{-1}AL, \quad B' = L^{-1}BL,
\]

so folgt aus den Gesetzen der Composition:

\[
A'B' = L^{-1}ABL,
\]

und daraus also der Satz:

\par
7.\enspace\gesperrt{Durchläuft $A$ die Substitutionen einer Gruppe, so durchläuft bei feststehendem $L$ die Transformirte $L^{-1}AL$ die Substitutionen einer isomorphen Gruppe.}

Von den inversen Substitutionen sind wohl zu unterscheiden die \gesperrt{transponirten Substitutionen}.

Man erhält nämlich aus jeder Substitution $A$ eine bestimmte andere, die die transponirte Substitution zu $A$ heisst, und die wir für den Augenblick mit $A_1$ bezeichnen wollen, wenn man in $A$ die Zeilen zu Colonnen macht, und umgekehrt, wenn man also in (2) die oberen mit den unteren Indices der $a$ vertauscht.

Wenn man in der Summe (11), $\sum^i b_i^{(k)} a_h^{(i)}$, die unteren mit den oberen Indices und gleichzeitig $a$ mit $b$ vertauscht, so erhält man einen Ausdruck, der nach (11) gleich $c_k^{(h)}$ ist.
